\section{Введение}

Слабые магнитные поля, энергия которых в расчёте на атом ничтожна по
сравнению с тепловой энергией $kT$, тем не менее изменяют пластичность
немагнитных кристаллов; совокупность таких наблюдений известна как
магнитопластический эффект [1,\,2]. Алюминиевые сплавы с железосодержащими
включениями занимают в этом ряду особое место: если такое включение
ферромагнитно, они допускают на первый взгляд простой механический путь от
поля к пластическому течению --- поле намагничивает включение,
магнитострикция (изменение формы намагниченного тела) деформирует его, и
напряжение, вызванное несоответствием между деформированным включением и
окружающей матрицей, передаётся матрице. Вопрос о том, действительно ли
включения в этих сплавах ферромагнитны, рассматривается в разделе~5.1. Серия экспериментов на промышленном сплаве Al--Mg--Si с
железосодержащими микровключениями подробно описала явление: 30-минутная
выдержка в поле $B = 0{,}4$--$0{,}7$~Тл, проводимая \emph{до} механических
испытаний, увеличивает последующую деформацию ползучести (медленного
деформирования под постоянной нагрузкой) при 160~МПа и комнатной температуре
примерно на четверть, в несколько раз повышает тепловыделение при
деформировании и смещает коэффициент Тейлора--Куинни --- долю работы
пластической деформации, выделяющуюся в виде тепла, --- с 0{,}07 до 0{,}25
[3--6]. Поскольку эффект сохраняется после снятия поля, его называют
«магнитной памятью» пластичности [4].

Предложенная в этих работах интерпретация опирается на единственную оценку.
Напряжение на границе включения с матрицей, создаваемое магнитострикцией
включения и полученное из величины поля через эмпирический коэффициент
пропорциональности, оценено примерно в 147~МПа при $B = 0{,}7$~Тл [5]. Это
превышает предел текучести матрицы 120~МПа, откуда и сделан вывод, что
матрица вокруг каждого включения пластически деформируется, а плотность
дислокаций --- линейных дефектов кристаллической решётки, движение которых и
составляет пластическое течение, --- возрастает. Эта цепочка рассуждений
никогда не проверялась на том масштабе, на котором она должна действовать.
Возникает ли в матрице напряжение такой величины, на каком расстоянии от
границы оно сохраняется и способно ли оно зарождать дислокации или приводить
в движение уже имеющиеся --- всё это вопросы об ангстремах и пикосекундах,
то есть о естественной области классической молекулярной динамики (МД), в
которой атомы движутся по уравнениям Ньютона под действием сил, задаваемых
межатомным потенциалом.

В настоящей работе магнитострикционный механизм подвергнут количественной
атомистической проверке на границе между алюминием и интерметаллидом
Al$_{13}$Fe$_4$. Магнитное поле в классической МД непосредственно не
представлено, поэтому его действие задаётся механически: по данным более
ранней экспериментальной работы [5] включение считается удлинённым вдоль
направления поля на $\varepsilon = 1{,}94\cdot10^{-3}$ (0{,}194\%) при
неизменном объёме --- это деформация, отвечающая оценённому там напряжению
на границе 147~МПа. Прямое воспроизведение макроскопического эффекта в
расчётной ячейке не предпринимается; причина изложена в разделе
«Обсуждение». Вместо этого каждое звено предполагаемой цепочки исследуется в
отдельности. Рассчитываются: поле напряжений, которое удлинённое включение
создаёт в матрице, и его затухание с расстоянием от границы; приложенное
сдвиговое напряжение, при котором уже имеющаяся в матрице дислокация
приходит в движение, и образуются ли на границе новые дислокации вплоть до наибольшего приложенного напряжения; сопротивление, которое растворённые в алюминии атомы Mg и Si
оказывают движению дислокации. Затем цепочка замыкается с противоположного
конца: двухмасштабная оценка даёт напряжение на границе, которое
\emph{требуется} для измеренного прироста ползучести на 25\% при термически
активированном скольжении --- движении дислокаций, при котором препятствия
преодолеваются с помощью тепловых флуктуаций, так что скорость скольжения
экспоненциально зависит от напряжения, --- и при затухании упругого поля
вне деформированного включения как $r^{-3}$ [7].